% Nouveau EGA V (cote 25): archive cote 25, site page 134 (PDF page 135).
% English translation of the checked French annotated master.
% AI-translated by GPT-6 Astra (OpenAI, ChatGPT Pro mode).
\scanpage{25/134}
\mspage{\hand{37.}}

\var{\margin{29.9.1965}}{\placedmargin{29.9.1965}}{marginal note placed}{ega25_25_134-v0001}

Dear Dieudonné,

Thank you for your letter of the 24\textsuperscript{th} and for the table of contents of sections
16 to 19. I would be glad to receive, when you have an opportunity, the provisional table
of contents of sections 20 and 21; I agree to include them in
fascicle 4 of Chap IV. But how do you want to subdivide \var{\xout{them}}{}{typewriter strike-out}{ega25_25_134-v0003}\var{\inserted[typed]{my old section 20}}{ my old section 20}{insertion}{ega25_25_134-v0004}, and
what will the titles of the two parts be \unsure{?|2} Since I am beginning to
lose track of the plan, and it is sometimes convenient to \unsure{be able to} refer
to a section no. without screwing up too much, I give you here what
seems to me to be the current plan; tell me whether you agree:

20. ???

21. ???

22. Linear systems, complements on the Picard group.

23. Grassmannians.

24. Smooth forms, ordinary quadratic singularities.

25. Hyperplane sections and all that crap.

26. Resultant and discriminant.

27. Infinitesimal extensions.

Section 25, \unsure{incidentally}, may well be very long, and I can already see you
wanting to subdivide it into \unsure{two}! Yet $27=3^{\hand{3}}$ is such a pretty
number!

There is no question of my publishing the former Appendix to section 18
under my name; your exposition has almost nothing left in common with
the vague handwritten notes I had given you, if indeed I ever
gave you any, and did not merely say to you: you only have to do

\mspage{\hand{134}}
\tnote{25-134-1}{A canceled doubtful d before d'accord is a French letter-level cancellation with no English counterpart. The live uncertainties rendered as ``be able to'', ``incidentally'', and ``two'' correspond respectively to pauvoir/pouvoir, the a/q within d'ailleurs, and the final eux in deux. The source also has the spellings paffois, Systêmes, manuscriptes, and j-amais; these have not been imitated with invented English misspellings. The informal wording ``without screwing up too much'' and ``all that crap'' renders sans trop déconner and bordel.}
